% GENERATED FILE. Edit canonical lessons, then run npm run generate:books.
% canonical-source-hash: fnv1a64:21d88ccc840f5274
% canonical-lessons: FA-C353-dowran, FA-C353-sade, FA-C353-movarrekh, FA-C353-niyakan, FA-C353-miras

\chapter{Eras and ancestors}
\label{ch:fa-eras-and-ancestors}
\hlchaptermodality{353}

\begin{chapteropening}
\textbf{By the end of this chapter:} \emph{I can talk about an era, a century, a historian, ancestors and heritage.}

Complete the last lesson of chapter 353: I can talk about an era, a century, a historian, ancestors and heritage.
\end{chapteropening}

\section[\texorpdfstring{\fa{دوران} (dowrân)}{دوران (dowrân)}]{\fa{دوران} (dowrân) — an era, a period}

\label{lesson:FA-C353-dowran}



\begin{quote}
Before the new one: say the Persian for an empire, then the Persian for a dynasty.
\end{quote}

\subsection*{You'll want to know: \fa{دوران}}
\textbf{\fa{دوران}} (\emph{dowrân}) — \textquotedblleft{}an era, a period\textquotedblright{}.

\begin{quote}
\fa{در} \fa{دوران} \fa{جنگ} \fa{زندگی} \fa{سخت} \fa{بود}.
\end{quote}

\begin{quote}
\emph{Dar dowrân-e jang zendegi sakht bud.}
\end{quote}

\begin{quote}
During the war years life was hard.
\end{quote}

\emph{Dowrân-e kudaki} \textquotedblleft{}childhood\textquotedblright{}.

\begin{cousinweb}
\emph{dowrân} is usually traced to Arabic \emph{dawarān}, \textquotedblleft{}a turning, a cycle\textquotedblright{}; EtymDB lists no entry.
\end{cousinweb}

\begin{grammarlens}[title={What you've built}]
One more word for telling, explaining and describing at B1.
\end{grammarlens}

\subsection*{Your turn}
\begin{itemize}
\raggedright
  \item \emph{Say it:} \emph{dowrân}
  \item \emph{Say it:} \emph{dowrân}, once more
  \item \emph{Say it:} a sentence of your own with \emph{dowrân}
  \item \emph{Recall:} say the Persian for an empire, then the Persian for a dynasty, then say \emph{dowrân} again
\end{itemize}

\begin{tcolorbox}[breakable,colback=teal!4,colframe=teal!35!black,title={Before you move on}]
What does \fa{دوران} mean? (\textquotedblleft{}An era, a period\textquotedblright{}.) Say it once more.
\end{tcolorbox}
\section[\texorpdfstring{\fa{سده} (sade)}{سده (sade)}]{\fa{سده} (sade) — a century}

\label{lesson:FA-C353-sade}



\begin{quote}
Before the new one: say the Persian for a dynasty, then the Persian for an era.
\end{quote}

\subsection*{You'll want to know: \fa{سده}}
\textbf{\fa{سده}} (\emph{sade}) — \textquotedblleft{}a century\textquotedblright{}.

\begin{quote}
\fa{این} \fa{مسجد} \fa{در} \fa{سده} \fa{دوازدهم} \fa{ساخته} \fa{شد}.
\end{quote}

\begin{quote}
\emph{In masjed dar sade-ye davâzdahom sâkhte shod.}
\end{quote}

\begin{quote}
This mosque was built in the twelfth century.
\end{quote}

The native word for a century; \emph{qarn} is the Arabic one. Also the name of a winter fire festival.

\begin{cousinweb}
\emph{sade} is a native Persian word that EtymDB lists without a parent; it is usually linked with \emph{sad} \textquotedblleft{}hundred\textquotedblright{}.
\end{cousinweb}

\begin{grammarlens}[title={What you've built}]
One more word for telling, explaining and describing at B1.
\end{grammarlens}

\subsection*{Your turn}
\begin{itemize}
\raggedright
  \item \emph{Say it:} \emph{sade}
  \item \emph{Say it:} \emph{sade}, once more
  \item \emph{Say it:} a sentence of your own with \emph{sade}
  \item \emph{Recall:} say the Persian for a dynasty, then the Persian for an era, then say \emph{sade} again
\end{itemize}

\begin{tcolorbox}[breakable,colback=teal!4,colframe=teal!35!black,title={Before you move on}]
What does \fa{سده} mean? (\textquotedblleft{}A century\textquotedblright{}.) Say it once more.
\end{tcolorbox}
\section[\texorpdfstring{\fa{مورخ} (movarrekh)}{مورخ (movarrekh)}]{\fa{مورخ} (movarrekh) — a historian}

\label{lesson:FA-C353-movarrekh}



\begin{quote}
Before the new one: say the Persian for an era, then the Persian for a century.
\end{quote}

\subsection*{You'll want to know: \fa{مورخ}}
\textbf{\fa{مورخ}} (\emph{movarrekh}) — \textquotedblleft{}a historian\textquotedblright{}.

\begin{quote}
\fa{مورخ} \fa{ها} \fa{درباره} \fa{این} \fa{جنگ} \fa{زیاد} \fa{نوشته} \fa{اند}.
\end{quote}

\begin{quote}
\emph{Movarrekh-hâ darbâre-ye in jang ziyâd neveshte-and.}
\end{quote}

\begin{quote}
Historians have written a lot about this war.
\end{quote}

Also plural \emph{movarrekhân}. The same Arabic root gives \emph{târikh} \textquotedblleft{}history\textquotedblright{}.

\begin{cousinweb}
\emph{movarrekh} is usually traced to Arabic \emph{mu'arrikh}, \textquotedblleft{}historian\textquotedblright{}, of the same root as \emph{ta'rīkh} \textquotedblleft{}history\textquotedblright{}; EtymDB lists no entry.
\end{cousinweb}

\begin{grammarlens}[title={What you've built}]
One more word for telling, explaining and describing at B1.
\end{grammarlens}

\subsection*{Your turn}
\begin{itemize}
\raggedright
  \item \emph{Say it:} \emph{movarrekh}
  \item \emph{Say it:} \emph{movarrekh}, once more
  \item \emph{Say it:} a sentence of your own with \emph{movarrekh}
  \item \emph{Recall:} say the Persian for an era, then the Persian for a century, then say \emph{movarrekh} again
\end{itemize}

\begin{tcolorbox}[breakable,colback=teal!4,colframe=teal!35!black,title={Before you move on}]
What does \fa{مورخ} mean? (\textquotedblleft{}A historian\textquotedblright{}.) Say it once more.
\end{tcolorbox}
\section[\texorpdfstring{\fa{نیاکان} (niyâkân)}{نیاکان (niyâkân)}]{\fa{نیاکان} (niyâkân) — ancestors, forefathers}

\label{lesson:FA-C353-niyakan}



\begin{quote}
Before the new one: say the Persian for a century, then the Persian for a historian.
\end{quote}

\subsection*{You'll want to know: \fa{نیاکان}}
\textbf{\fa{نیاکان}} (\emph{niyâkân}) — \textquotedblleft{}ancestors, forefathers\textquotedblright{}.

\begin{quote}
\fa{نیاکان} \fa{ما} \fa{در} \fa{این} \fa{روستا} \fa{زندگی} \fa{میکردند}.
\end{quote}

\begin{quote}
\emph{Niyâkân-e mâ dar in rustâ zendegi mikardand.}
\end{quote}

\begin{quote}
Our ancestors lived in this village.
\end{quote}

The plural of the rare \emph{niyâ} \textquotedblleft{}grandfather\textquotedblright{}. Literary.

\begin{cousinweb}
\emph{niyâkân} is the plural of \emph{niyâ} \textquotedblleft{}grandfather, forebear\textquotedblright{}, a native Persian word; EtymDB lists no entry.
\end{cousinweb}

\begin{grammarlens}[title={What you've built}]
One more word for telling, explaining and describing at B1.
\end{grammarlens}

\subsection*{Your turn}
\begin{itemize}
\raggedright
  \item \emph{Say it:} \emph{niyâkân}
  \item \emph{Say it:} \emph{niyâkân}, once more
  \item \emph{Say it:} a sentence of your own with \emph{niyâkân}
  \item \emph{Recall:} say the Persian for a century, then the Persian for a historian, then say \emph{niyâkân} again
\end{itemize}

\begin{tcolorbox}[breakable,colback=teal!4,colframe=teal!35!black,title={Before you move on}]
What does \fa{نیاکان} mean? (\textquotedblleft{}Ancestors, forefathers\textquotedblright{}.) Say it once more.
\end{tcolorbox}
\section[\texorpdfstring{\fa{میراث} (mirâs)}{میراث (mirâs)}]{\fa{میراث} (mirâs) — heritage, a legacy}

\label{lesson:FA-C353-miras}



\begin{quote}
Before the new one: say the Persian for a historian, then the Persian for ancestors.
\end{quote}

\subsection*{You'll want to know: \fa{میراث}}
\textbf{\fa{میراث}} (\emph{mirâs}) — \textquotedblleft{}heritage, a legacy\textquotedblright{}.

\begin{quote}
\fa{این} \fa{بنای} \fa{قدیمی} \fa{میراث} \fa{همه} \fa{مردم} \fa{است}.
\end{quote}

\begin{quote}
\emph{In banâ-ye qadimi mirâs-e hame-ye mardom ast.}
\end{quote}

\begin{quote}
This old building is the heritage of all the people.
\end{quote}

Written with \fa{ث}. \emph{Mirâs-e farhangi} \textquotedblleft{}cultural heritage\textquotedblright{}.

\begin{cousinweb}
\emph{mirâs} is borrowed from Arabic \emph{mīrāth}, \textquotedblleft{}inheritance\textquotedblright{} (EtymDB).
\end{cousinweb}

\begin{grammarlens}[title={What you've built}]
One more word for telling, explaining and describing at B1.
\end{grammarlens}

\subsection*{Your turn}
\begin{itemize}
\raggedright
  \item \emph{Say it:} \emph{mirâs}
  \item \emph{Say it:} \emph{mirâs}, once more
  \item \emph{Say it:} a sentence of your own with \emph{mirâs}
  \item \emph{Recall:} say the Persian for a historian, then the Persian for ancestors, then say \emph{mirâs} again
\end{itemize}

\begin{tcolorbox}[breakable,colback=teal!4,colframe=teal!35!black,title={Before you move on}]
What does \fa{میراث} mean? (\textquotedblleft{}Heritage, a legacy\textquotedblright{}.) Say it once more.
\end{tcolorbox}
